\documentclass{YM_Likhanova}
\usepackage{amssymb,amsmath,bm}
\usepackage{cmap}
\usepackage{cite}

\setenumerate{topsep=2pt,itemsep=2pt,parsep=0pt,labelindent=\parindent,leftmargin=*}
\setlist{topsep=2pt,itemsep=2pt,parsep=0pt,labelindent=\parindent,leftmargin=*}

\newcommand{\const}{{\rm const\,}}
\newcommand{\grad}{{\rm grad\,}}
\renewcommand{\vec}[1]{\boldsymbol{#1}}
\newcommand{\ds}{\displaystyle}
\newcommand{\No}{№}

\THeader{\small{XXI Всероссийская конференция молодых учёных по математическому моделированию \linebreak и информационным технологиям. Новосибирск, 7 -- 11 декабря 2020 г.}\\[-2mm]}

\begin{document}

\Authors{Автор1 И.О., Автор2 И.О.}
\Title{Название доклада}
\Text{
% в систему Конференции текст вставлять НАЧИНАЯ С этого места!
После редактирования тезисов вставьте основную часть (то, что в поле \textbackslash{}Text\{\}) в соответствующее поле системы Конференции, убедившись, что стоит галочка <<в формате LaTex>>. После сохранения {\bf обязательно проверьте создавшийся в системе PDF!}

В тексте тезисов {\bf необходимо} кратко описать, какая задача решалась, какими методами и какие результаты были получены. {\it Желательно} отметить актуальность решения поставленной задачи и научную новизну выполненной работы.

Рекомендуемый объём тезисов вместе с названием и литературой~--- один столбец.
Максимальный~--- полтора столбца. Рисунки не допускаются.

Пример оформления нумерованной формулы:
\begin{equation}
\label{eq_SvoyaFamiliyaInicialy_1}
  {\bf A x} = {\bf b},
\end{equation}
и ссылки на неё~(\ref{eq_SvoyaFamiliyaInicialy_1}). 
Если на формулу нет ссылки в тексте, она должна быть ненумерованной:
\begin{equation*}
\frac{\partial y}{\partial t} = -\frac{1}{r}\sqrt{1+\left(\frac{\partial h}{\partial r}\right)^2},\quad \mbox{где } h = g(r).
\end{equation*}

Пример нумерованного списка:
\begin{enumerate}
\item Текст пункта 1.
\item Текст пункта 2.
\end{enumerate}

Пример ненумерованного списка:
\begin{itemize}
\item текст пункта 1;
\item текст пункта 2.
\end{itemize}

Составные названия уравнений и т.д. пишутся через тире с пробелами: уравнения Навье~--- Стокса.

{ \it Список литературы} должен состоять не более чем из 4 ссылок.
Примеры оформления: статья на русском~\cite{ct_SvoyaFamiliyaInicialy_1} и 
на английском \cite{ct_SvoyaFamiliyaInicialy_2},
монография на русском	(книга)~\cite{ct_SvoyaFamiliyaInicialy_3} и
на английском~\cite{ct_SvoyaFamiliyaInicialy_4},
электронный ресурс~\cite{ct_SvoyaFamiliyaInicialy_5},
диссертация на русском~\cite{ct_SvoyaFamiliyaInicialy_6} и авторефереат~\cite{ct_SvoyaFamiliyaInicialy_7},
диссертация на английском~\cite{ct_SvoyaFamiliyaInicialy_8},
труды конференции на русском~\cite{ct_SvoyaFamiliyaInicialy_9} и на английском~\cite{ct_SvoyaFamiliyaInicialy_10},
статья или глава в книге под редакцией~\cite{ct_SvoyaFamiliyaInicialy_11},
книга под редакцией~\cite{ct_SvoyaFamiliyaInicialy_12},
свидетельство об интеллектуальной собственности~\cite{ct_SvoyaFamiliyaInicialy_13},
ГОСТ~\cite{ct_SvoyaFamiliyaInicialy_14}. Если количество авторов превышает четыре, то пишутся первые три и следом <<и~др.>> (<<et~al.>> в англоязычной ссылке). 
% ct_SvoyaFamiliyaInicialy - уникальный идентификатор для каждого докалада!
% Если у Вас несколько докладов на этой конференции, то необходимо указать номер доклада в идентификаторе. Например так: ct_Nomer_SvoyaFamiliyaInicialy.

Оформление ссылки на грант:

\emph{Работа выполнена при финансовой поддержке РНФ (грант \No~ХХ-ХХ-XXXXX).}

Научного руководителя рекомендуется не включать в соавторы, а отмечать перед списком литературы:

\emph{Научный руководитель~--- д.ф.-м.н. Фамилия~И.\,О.}

\begin{thebibliography}{9}
%Статья на русском:
\bibitem{ct_SvoyaFamiliyaInicialy_1}
{\sc ФамилияАвтора1~И.\,О., ФамилияАвтора2~И.\,О.}
Название статьи~// Название журнала. 2014. Т.~17. №~4~(60). С.~38--47.
 
%Статья на английском:
\bibitem{ct_SvoyaFamiliyaInicialy_2}
{\sc AuthorSurname1~N., AuthorSurname2~N.}
Article title~// Journal name. 2017. Vol.~10. N.~1. P.~26--46.

%Монография на русском	(книга):
\bibitem{ct_SvoyaFamiliyaInicialy_3}
{\sc ФамилияАвтора1~И.\,О., ФамилияАвтора2~И.\,О.}
Название монографии~/ Город:~Издательство, 2015. 204~с.
 
%Монография на английском:
\bibitem{ct_SvoyaFamiliyaInicialy_4}
{\sc AuthorSurname~N.\,N.}
Monograph title~/ Location: Publisher, 1992. 259~p.
 
%Электронный ресурс:
\bibitem{ct_SvoyaFamiliyaInicialy_5}
{\sc ФамилияАвтора~И.\,О.}
Название. [Электронный ресурс]. URL: http://adress.html (дата обращения 01.02.2018).
 
%Диссертация на русском и авторефереат:
\bibitem{ct_SvoyaFamiliyaInicialy_6}
{\sc ФамилияАвтора~И.\,О.}
Название работы: дис. ... канд./д-ра физ.-мат. наук. В какой организации выполнялась, Город, 2005. 130~с.

\bibitem{ct_SvoyaFamiliyaInicialy_7}
{\sc ФамилияАвтора~И.\,О.}
Название работы: автореф. дис. ... канд./д-ра физ.-мат. наук. В какой организации выполнялась, Город, 2005. 15~с.

%Диссертация на английском:
\bibitem{ct_SvoyaFamiliyaInicialy_8}
{\sc AuthorSurname N.}
Title. PhD theses. Name of Institution, Location, Year. 150~p.
 
%Труды конференции на русском:
\bibitem{ct_SvoyaFamiliyaInicialy_9}
{\sc ФамилияАвтора1~И.\,О., ФамилияАвтора2~И.\,О., ФамилияАвтора3~И.\,О. и др.}
Название доклада~// Тр. V~Междунар. Конф. <<Название конференции>>. Город:~Издательство, 2013. Т.~2. С.~45--55.
 
%Труды конференции на английском:
\bibitem{ct_SvoyaFamiliyaInicialy_10}
{\sc AuthorSurname~N.\,N.}
Proceedings title~// Proc. Intern. Conf. <<Conference name>>. Location: Publisher, 2018. Vol.~1. N.~5. P.~15--22.

%Статья или глава в книге под редакцией:
\bibitem{ct_SvoyaFamiliyaInicialy_11}
{\sc AuthorSurname~N.\,N.}
Title of chapter~/ In: N.\,N.~EditorSurname1, N.\,N.~EditorSurname2 (Eds.) Title of book: Subtitle. Location:~Publisher, 2012. P.~33--135.

%Книга под редакцией:
\bibitem{ct_SvoyaFamiliyaInicialy_12}
Название Книги~/ Под ред. И.\,О.~ФамилияРедактора. Город:~Издательство, 1987. 712~с. 

%Cвидетельство интеллектуальной собственности:
\bibitem{ct_SvoyaFamiliyaInicialy_13}
{\sc ФамилияАвтора~И.\,О.}
Название собственности (свидетельство \No~2017615165)~/ М.:~Федеральная служба по интеллектуальной собственности (Роспатент), 2017.

%ГОСТ
\bibitem{ct_SvoyaFamiliyaInicialy_14}
ГОСТ Р~22.1.06-99. Безопасность в чрезвычайных ситуациях. Мониторинг и прогнозирование опасных геологических явлений и процессов. Общие требования. Введён~2000-01-01~/ М.:~Госстандарт России: Изд-во стандартов, 1999. 25~с. 

\end{thebibliography}
% в систему Конференции текст вставлять ДО этого места!
} % конец \Text
\PutItem
\EndDoc
\end{document}
